\chapter{Cosmic Energy Budget and Macroscopic Chronology}
\label{app:cosmic_energy_budget_and_chronology}
\addtocontents{toc}{\protect\setcounter{tocdepth}{0}}

\noindent\textbf{Abstract.}
This appendix develops the macroscopic geometric limits and structural
chronology of the Tension Field Theory (TFT) cosmic oscillator. Modeling the
active Spacetime Gravitational Floor (SGF) as a dual-phase elastodynamic
Face-Centered Cubic (FCC) medium undergoing localized FCC$\rightarrow$SC phase
transitions, we derive the local Octahedral accommodation limit, construct the
global mass-weighted percolation threshold for Catastrophic Decompression, and
formulate the Global Phase Avalanche from the integrated growth of the live
FCC/SC phase-boundary network.

The framework distinguishes the aggregate inertial mass carried by live
FCC/SC phase boundaries from the physical volume of the saturated SC cores.
The local live-boundary conservation law developed in \hyperref[app:live_black_hole_boundaries]{Appendix~\ref*{app:live_black_hole_boundaries}} leads to a
population-level structural growth rate
\(\Gamma_{\mathrm{pc}}\equiv\dot M_{\mathrm p}/M_{\mathrm p}\).
When the self-reproducing boundary network possesses an approximately
stationary net fractional growth eigenvalue, the aggregate phase-boundary mass
grows exponentially,
\[
    M_{\mathrm p}(t)
    =
    M_{\mathrm p}(0)
    e^{\Gamma_{\mathrm{pc}}t}.
\]

The percolation geometry fixes the structural avalanche depth required to
reach the Big Rupture, while \(\Gamma_{\mathrm{pc}}\) fixes the duration of one
structural e-fold. Under the single-rate Active-Vacuum synchronization
specialization,
\[
    \Gamma_{\mathrm{pc}}
    =
    \Gamma_{\mathrm D},
\]
the structural Phase Avalanche and coherent optical relaxation are governed
by the same macroscopic timescale without requiring direct conversion of
photon energy into Black Hole boundary mass.


% =====================================================================
\section{The Local Volumetric Yield Limit (Phase Condensation)}
% =====================================================================

As established in Part I, the active vacuum (State I) is an FCC lattice with
an invariant geometric packing fraction of
\begin{equation}
    \beta_{\mathrm{geo}}
    \approx
    0.7405.
\end{equation}

Mass formation requires localized kinematic excitation of the active lattice
and is supported by the interstitial degrees of freedom between the FCC
nodes. The primary isotropic accommodation pockets are the Octahedral voids.

To determine the maximum local geometric accommodation associated with these
sites, we calculate the precise volumetric fraction of the inscribed
Octahedral regions relative to the FCC unit cell.

\begin{itemize}

    \item \textbf{Deriving the Octahedral Void Radius
    (\(r_{\mathrm O}\)):}

    In an FCC lattice formed by nodes of radius \(r_{\mathrm D}\), neighboring
    nodes touch along the face diagonal. The face diagonal is therefore
    \(4r_{\mathrm D}\). By the Pythagorean theorem, the edge length of the
    conventional unit cell is
    \begin{equation}
        a
        =
        \sqrt{
            \frac{
                (4r_{\mathrm D})^2
            }{
                2
            }
        }
        =
        2\sqrt{2}\,r_{\mathrm D}.
    \end{equation}

    One class of FCC Octahedral sites lies at the midpoint of the unit-cell
    edges, with the symmetry-equivalent body-center site completing the
    Octahedral-site population of the conventional cell. Along an edge, the
    distance from the center of a corner node to the center of the Octahedral
    site is \(a/2\). Therefore,
    \begin{equation}
        r_{\mathrm D}
        +
        r_{\mathrm O}
        =
        \frac{a}{2}
        =
        \sqrt{2}\,r_{\mathrm D}.
    \end{equation}

    Isolating the inscribed radius gives
    \begin{equation}
        r_{\mathrm O}
        =
        (\sqrt{2}-1)
        r_{\mathrm D}.
    \end{equation}

    \item \textbf{Deriving the Volumetric Fraction:}

    The total conventional-cell volume is
    \begin{equation}
        V_{\mathrm{cell}}
        =
        a^3
        =
        16\sqrt{2}\,
        r_{\mathrm D}^3.
    \end{equation}

    There are four Octahedral sites per conventional FCC unit cell: one
    body-center site and twelve edge sites, each shared by four neighboring
    cells. The total inscribed Octahedral volume is therefore
    \begin{equation}
        V_{\mathrm{oct}}
        =
        4
        \left(
            \frac{4}{3}
            \pi
            r_{\mathrm O}^3
        \right)
        =
        \frac{16\pi}{3}
        (\sqrt{2}-1)^3
        r_{\mathrm D}^3.
    \end{equation}

    Dividing by the unit-cell volume gives the Octahedral Volumetric Fraction,
    \begin{equation}
        \beta_{\mathrm{oct}}
        =
        \frac{
            V_{\mathrm{oct}}
        }{
            V_{\mathrm{cell}}
        }
        =
        \frac{
            \pi(\sqrt{2}-1)^3
        }{
            3\sqrt{2}
        }
        \approx
        0.0526.
    \end{equation}

\end{itemize}

Within the TFT structural-capacity interpretation, this \(5.26\%\) fraction
provides the local isotropic Octahedral accommodation ceiling,
\begin{equation}
    \beta_{\mathrm b}
    \leq
    \beta_{\mathrm{oct}}
    \approx
    0.0526.
\end{equation}

This quantity is a local geometric capacity benchmark. It is not the physical
volume fraction of saturated State II SC cores. Black Hole formation occurs
when the local active FCC state reaches gravitational saturation,
\begin{equation}
    \alpha_{\mathrm G}
    \rightarrow
    1,
\end{equation}
forcing the FCC$\rightarrow$SC phase transition and creating a live FCC/SC
phase boundary.


% =====================================================================
\section{The Global Topological Limit
(Percolation Theory and The Big Rupture)}
% =====================================================================

The local Octahedral limit constrains the geometric accommodation of localized
matter defects. The ultimate global stability of the macroscopic reservoir is
instead governed by connectivity of the distributed State II phase-boundary
network.

In TFT, the inertial mass of a Black Hole is phase-segregated onto its live
FCC/SC boundary. The global quantity relevant to percolation is therefore the
Aggregate Percolating Mass \(M_{\mathrm p}\), defined as the combined inertial
mass carried by the participating live phase boundaries:
\begin{equation}
    M_{\mathrm p}(t)
    \equiv
    \sum_i
    \int_{\partial\Omega_i(t)}
    \Sigma_{\partial}
    \,\mathrm dA,
\end{equation}
where \(\Sigma_{\partial}\) is the local surface mass density and the sum is
taken over the distinct live FCC/SC boundary domains.

Once initially separate boundaries merge through a saturated bridge, the
resulting common boundary is treated as one connected phase-boundary domain.

For separated approximately spherical Black Holes, this reduces to
\begin{equation}
    M_{\mathrm p}(t)
    =
    \sum_i
    M_i(t).
\end{equation}

Thus \(M_{\mathrm p}\) is an aggregate boundary-mass measure. It is not the
summed geometrical volume of the SC interiors.

\begin{itemize}

    \item \textbf{The Interwoven Interstitial Network:}

    The interstitial regions of the FCC Soft Floor form a connected geometric
    network through which localized phase saturation and boundary coupling can
    develop. State II cores nucleate locally, while their exterior strain
    fields and live phase boundaries can subsequently couple across this
    network.

    \item \textbf{Site Percolation Threshold (\(p_{\mathrm c}\)):}

    TFT adopts the FCC site-percolation value
    \begin{equation}
        p_{\mathrm c}
        \approx
        0.199
    \end{equation}
    as the baseline coarse-grained connectivity threshold for the distributed
    defect network.

    Because the actual cosmological phase domains grow and interact through
    correlated strain fields, a complete nonlinear correlated-percolation
    calculation remains a future microscopic refinement of this baseline
    closure.

    \item \textbf{Total Interstitial Fraction (\(\beta_{\mathrm{IC}}\)):}

    The total geometrical interstitial fraction of the active FCC vacuum is
    \begin{equation}
        \beta_{\mathrm{IC}}
        =
        1-\beta_{\mathrm{geo}}
        \approx
        0.2595.
    \end{equation}

\end{itemize}

To connect the site-occupation probability to the macroscopic boundary-mass
fraction, TFT adopts an equal-capacity coarse-graining of the interstitial
network.

Let \(N_{\mathrm{cap}}\) denote the number of equivalent coarse-grained
interstitial capacity elements and let \(m_{\mathrm{cap}}\) denote the
structural mass-equivalent capacity assigned to each element. Then
\begin{equation}
    \beta_{\mathrm{IC}}
    M_{\mathrm{Res}}
    =
    N_{\mathrm{cap}}
    m_{\mathrm{cap}}.
\end{equation}

If \(N_{\mathrm{occ}}\) of these elements participate in the connected State
II boundary network, their aggregate boundary mass is
\begin{equation}
    M_{\mathrm p}
    =
    N_{\mathrm{occ}}
    m_{\mathrm{cap}}.
\end{equation}

Dividing these relations gives
\begin{equation}
    \frac{
        M_{\mathrm p}
    }{
        \beta_{\mathrm{IC}}
        M_{\mathrm{Res}}
    }
    =
    \frac{
        N_{\mathrm{occ}}
    }{
        N_{\mathrm{cap}}
    }.
\end{equation}

Defining the coarse-grained occupation fraction
\begin{equation}
    p
    \equiv
    \frac{
        N_{\mathrm{occ}}
    }{
        N_{\mathrm{cap}}
    },
\end{equation}
we therefore obtain
\begin{equation}
    p
    =
    \frac{
        M_{\mathrm p}
    }{
        \beta_{\mathrm{IC}}
        M_{\mathrm{Res}}
    }.
\end{equation}

Under this equal-capacity FCC percolation closure, global connectivity occurs
when
\begin{equation}
    p
    =
    p_{\mathrm c}.
\end{equation}

The corresponding Global Rupture Trigger Mass is therefore
\begin{equation}
    \boxed{
    M_{\mathrm{trg}}
    =
    p_{\mathrm c}
    \beta_{\mathrm{IC}}
    M_{\mathrm{Res}}.
    }
\end{equation}

Defining the rupture fraction relative to the total reservoir mass,
\begin{equation}
    \beta_{\mathrm{rup}}
    \equiv
    \frac{
        M_{\mathrm{trg}}
    }{
        M_{\mathrm{Res}}
    },
\end{equation}
gives
\begin{equation}
    \beta_{\mathrm{rup}}
    =
    p_{\mathrm c}
    \beta_{\mathrm{IC}}
    =
    0.199
    \times
    0.2595
    \approx
    0.0516.
\end{equation}

The physical meaning of this \(5.16\%\) threshold is therefore
\begin{equation}
    \frac{
        M_{\mathrm p}
    }{
        M_{\mathrm{Res}}
    }
    \rightarrow
    0.0516,
\end{equation}
not
\begin{equation}
    \frac{
        V_{\mathrm{SC}}
    }{
        V_{\mathrm{Res}}
    }
    \rightarrow
    0.0516.
\end{equation}

The global rupture threshold is an aggregate live-boundary mass threshold
expressed in the common structural-capacity normalization of the reservoir.

Within that same structural-capacity normalization,
\begin{equation}
    \beta_{\mathrm{rup}}
    \approx
    0.0516
    <
    \beta_{\mathrm{oct}}
    \approx
    0.0526.
\end{equation}

Thus, under the baseline equal-capacity percolation closure, the global
connectivity threshold lies slightly below the local Octahedral accommodation
ceiling: the distributed phase-boundary network reaches the global fracture
condition before the local structural-capacity benchmark is exhausted.


% ---------------------------------------------------------------------
\subsection{Derivation of the Deflation Parameters}
\label{subsec:deflation_parameters}
% ---------------------------------------------------------------------

The \(5.16\%\) rupture threshold is a boundary-mass fraction and must not be
interpreted as a \(5.16\%\) physical SC-core volume fraction.

The macroscopic deflation law instead follows from the adopted
structural-capacity ledger of the remaining active FCC reservoir:
\begin{equation}
    M_{\mathrm{FCC}}(t)
    =
    M_{\mathrm{Res}}
    -
    M_{\mathrm p}(t).
\end{equation}

For the quasi-static pre-rupture branch, TFT adopts the reservoir-density
closure that the remaining active FCC domain preserves its reference mean
structural density during global contraction:
\begin{equation}
    \frac{
        M_{\mathrm{FCC}}(t)
    }{
        V(t)
    }
    =
    \frac{
        M_{\mathrm{Res}}
    }{
        V_{\mathrm{Res}}
    }.
\end{equation}

It follows that
\begin{equation}
    \frac{
        V(t)
    }{
        V_{\mathrm{Res}}
    }
    =
    1
    -
    \frac{
        M_{\mathrm p}(t)
    }{
        M_{\mathrm{Res}}
    }.
\end{equation}

At the global rupture threshold,
\begin{equation}
    \frac{
        M_{\mathrm p}
    }{
        M_{\mathrm{Res}}
    }
    =
    \beta_{\mathrm{rup}}
    \approx
    0.0516,
\end{equation}
so
\begin{equation}
    V_{\mathrm{rup}}
    =
    (1-\beta_{\mathrm{rup}})
    V_{\mathrm{Res}}
    =
    0.9484
    V_{\mathrm{Res}}.
\end{equation}

Using the established reservoir volume,
\begin{equation}
    V_{\mathrm{rup}}
    \approx
    7.9058\times10^{104}
    \ {\rm m^3}.
\end{equation}

The corresponding spherical reservoir radius satisfies
\begin{equation}
    R_{\mathrm{rup}}
    =
    R_{\mathrm{Res}}
    (1-\beta_{\mathrm{rup}})^{1/3},
\end{equation}
giving
\begin{equation}
    R_{\mathrm{rup}}
    \approx
    5.7361\times10^{34}
    \ {\rm m}.
\end{equation}

The pre-rupture contraction of the macroscopic structural reservoir is
therefore
\begin{equation}
    \Delta R
    =
    R_{\mathrm{Res}}
    -
    R_{\mathrm{rup}}
    \approx
    1.0220\times10^{33}
    \ {\rm m}.
\end{equation}

These deflation values follow from the quasi-static mean-reservoir-density
closure of the remaining active FCC domain. They do not require a
\(5.16\%\) physical filling of the reservoir by zero-void SC cores.

The Global Rupture Trigger Mass itself is
\begin{equation}
    M_{\mathrm{trg}}
    =
    0.0516
    \left(
        1.6707\times10^{79}
        \ {\rm kg}
    \right)
    \approx
    8.6208\times10^{77}
    \ {\rm kg}.
\end{equation}


% ---------------------------------------------------------------------
\subsection{The Spanning Cluster and Topological Continuity}
% ---------------------------------------------------------------------

In percolation theory, the critical threshold marks the transition from
isolated clusters to a reservoir-spanning connected network. In TFT, this
connectivity does not require the physical SC interiors of individual Black
Holes to touch shoulder-to-shoulder.

The live-boundary merger mechanics instead permit two initially separated
State II domains to become mechanically connected when the strained FCC
region between them reaches saturation. The local bridge condition is
represented schematically by
\begin{equation}
    \alpha_{\mathrm{bridge}}
    =
    \alpha_{\mathrm{G},1}
    +
    \alpha_{\mathrm{G},2}
    +
    \alpha_{\mathrm{dyn}}
    \geq
    1,
\end{equation}
where
\(\alpha_{\mathrm{G},1}\)
and
\(\alpha_{\mathrm{G},2}\)
are the quasi-static gravitational strain contributions and
\(\alpha_{\mathrm{dyn}}\)
represents additional dynamical loading of the intervening region.

When this condition is satisfied across a connected path, the intermediate
active FCC region undergoes the same local saturation transition that defines
an individual Black Hole boundary:
\begin{equation}
    \mathrm{FCC}
    \rightarrow
    \mathrm{SC}.
\end{equation}

The global connectivity sequence is therefore
\begin{equation}
\begin{array}{c}
    \text{isolated State II domains}
    \\
    \Downarrow
    \\
    \text{overlapping exterior strain fields}
    \\
    \Downarrow
    \\
    \text{saturated FCC/SC bridges}
    \\
    \Downarrow
    \\
    \text{connected live-boundary network}
    \\
    \Downarrow
    \\
    \text{reservoir-spanning phase-boundary cluster}.
\end{array}
\end{equation}

The object that percolates is therefore the combined network of State II
cores, live FCC/SC boundaries, and saturated connecting bridges.

The physical SC interiors themselves need not occupy \(5.16\%\) of the
reservoir volume, nor must individual Black Hole cores be in direct geometric
contact.

The familiar term ``fault line'' is retained as an intuitive description of
this global fracture path. In the full three-dimensional reservoir, the
critical object is more properly a tortuous spanning fracture network or
manifold.


% ---------------------------------------------------------------------
\subsection{The Phase Trigger Unification}
% ---------------------------------------------------------------------

As the universe evolves, self-sustaining FCC$\rightarrow$SC phase domains
accumulate boundary mass and become increasingly interconnected.

Under the equal-capacity percolation closure, the baseline topological trigger
is reached when
\begin{equation}
    M_{\mathrm p}
    =
    M_{\mathrm{trg}},
\end{equation}
equivalently,
\begin{equation}
    \frac{
        M_{\mathrm p}
    }{
        \beta_{\mathrm{IC}}
        M_{\mathrm{Res}}
    }
    =
    p_{\mathrm c}.
\end{equation}

At this threshold the live-boundary and saturated-bridge system forms a
reservoir-spanning cluster.

TFT adopts the global rupture condition that the critical spanning cluster
becomes a load-severing fault network for the active FCC containment
structure. Consequently,
\begin{equation}
    \Delta P
    \rightarrow
    0.
\end{equation}

Loss of the stabilizing global pressure differential removes the mechanical
confinement of the phase-locked network. The hyper-compressed State II
structure then undergoes rapid martensitic snap-back,
\begin{equation}
    \mathrm{SC}
    \rightarrow
    \mathrm{FCC},
\end{equation}
producing Catastrophic Decompression and restoring the active State I
geometry for the next cosmic cycle.

The Big Rupture is therefore controlled by a mass-weighted topological
connectivity threshold, not by exhaustion of physical SC-core volume.


% =====================================================================
\section{The Chronology of the Global Phase Avalanche}
% =====================================================================

The percolation threshold determines the endpoint of the structural cycle.
To determine the corresponding macroscopic clock, TFT tracks the aggregate
live-boundary mass \(M_{\mathrm p}\) generated by the evolving
FCC$\rightarrow$SC phase-conversion network.


% ---------------------------------------------------------------------
\subsection{The Global Phase Avalanche Equation}
% ---------------------------------------------------------------------

Black Hole formation in TFT is a nucleated phase transition. Once a live
FCC/SC boundary has been established, its subsequent evolution is governed by
the local surface-mass conservation law developed in \hyperref[app:live_black_hole_boundaries]{Appendix~\ref*{app:live_black_hole_boundaries}}.

For one closed live boundary,
\begin{equation}
    M_{\partial}(t)
    =
    \int_{\partial\Omega(t)}
    \Sigma_{\partial}\,
    \mathrm dA.
\end{equation}

The integrated live-boundary conservation law gives
\begin{equation}
    \frac{\mathrm dM_{\partial}}{\mathrm dt}
    =
    \int_{\partial\Omega(t)}
    \left(
        \dot\Sigma_{\partial,\mathrm{in}}
        +
        \dot\Sigma_{\partial,\mathrm{vac}}
        -
        \dot\Sigma_{\partial,\mathrm{shed}}
    \right)
    \mathrm dA.
\end{equation}

For approximately spherical growth,
\begin{equation}
    R_{\mathrm{QG}}
    =
    \frac{
        GM_{\partial}
    }{
        c_{\mathrm E}^{2}
    },
\end{equation}
so the local fractional boundary-growth rate may be written
\begin{equation}
    \Gamma_{\partial}
    \equiv
    \frac{
        \dot M_{\partial}
    }{
        M_{\partial}
    }
    =
    \frac{
        \dot R_{\mathrm{QG}}
    }{
        R_{\mathrm{QG}}
    }
    =
    \frac{
        v_n
    }{
        R_{\mathrm{QG}}
    }.
\end{equation}

The cosmological phase avalanche is the population-level extension of this
live-boundary growth.

Let
\begin{equation}
    \mathcal J_{\mathrm{in}},
    \qquad
    \mathcal J_{\mathrm{vac}},
    \qquad
    \mathcal J_{\mathrm{shed}},
    \qquad
    \mathcal J_{\mathrm{nuc}}
\end{equation}
denote respectively the aggregate rates of boundary mass added by incoming
loading, FCC$\rightarrow$SC vacuum conversion, stress shedding, and creation
of newly participating State II nuclei.

The global structural balance is then
\begin{equation}
    \boxed{
    \frac{
        \mathrm dM_{\mathrm p}
    }{
        \mathrm dt
    }
    =
    \mathcal J_{\mathrm{in}}
    +
    \mathcal J_{\mathrm{vac}}
    -
    \mathcal J_{\mathrm{shed}}
    +
    \mathcal J_{\mathrm{nuc}}.
    }
    \label{eq:AU_phase_avalanche_master}
\end{equation}

Define the instantaneous Phase-Avalanche Rate by
\begin{equation}
    \boxed{
    \Gamma_{\mathrm{pc}}(t)
    \equiv
    \frac{1}{
        M_{\mathrm p}
    }
    \frac{
        \mathrm dM_{\mathrm p}
    }{
        \mathrm dt
    }.
    }
    \label{eq:AU_gamma_pc_definition}
\end{equation}

Thus,
\begin{equation}
    \boxed{
    \Gamma_{\mathrm{pc}}(t)
    =
    \frac{
        \mathcal J_{\mathrm{in}}
        +
        \mathcal J_{\mathrm{vac}}
        -
        \mathcal J_{\mathrm{shed}}
        +
        \mathcal J_{\mathrm{nuc}}
    }{
        M_{\mathrm p}
    }.
    }
    \label{eq:AU_gamma_pc_flux}
\end{equation}

Equation~\eqref{eq:AU_gamma_pc_flux}
is the structural definition of the macroscopic phase-avalanche rate.

In a statistically self-similar pre-rupture regime, the aggregate fluxes may
be represented by effective fractional rates,
\begin{equation}
    \mathcal J_{\mathrm{in}}
    =
    \Gamma_{\mathrm{in}}
    M_{\mathrm p},
\end{equation}
\begin{equation}
    \mathcal J_{\mathrm{vac}}
    =
    \Gamma_{\mathrm{vac}}
    M_{\mathrm p},
\end{equation}
\begin{equation}
    \mathcal J_{\mathrm{nuc}}
    =
    \Gamma_{\mathrm{nuc}}
    M_{\mathrm p},
\end{equation}
and
\begin{equation}
    \mathcal J_{\mathrm{shed}}
    =
    \Gamma_{\mathrm{shed}}
    M_{\mathrm p}.
\end{equation}

Then
\begin{equation}
    \boxed{
    \Gamma_{\mathrm{pc}}
    =
    \Gamma_{\mathrm{in}}
    +
    \Gamma_{\mathrm{vac}}
    +
    \Gamma_{\mathrm{nuc}}
    -
    \Gamma_{\mathrm{shed}}.
    }
    \label{eq:AU_gamma_pc_components}
\end{equation}

If this net reproduction eigenvalue is approximately stationary over the
long pre-rupture phase,
\begin{equation}
    \Gamma_{\mathrm{pc}}
    \simeq
    \text{constant},
\end{equation}
then
\begin{equation}
    \frac{
        \mathrm dM_{\mathrm p}
    }{
        \mathrm dt
    }
    =
    \Gamma_{\mathrm{pc}}
    M_{\mathrm p}.
\end{equation}

Integration gives the Global Phase Avalanche Equation,
\begin{equation}
    \boxed{
    M_{\mathrm p}(t)
    =
    M_{\mathrm p}(0)
    e^{\Gamma_{\mathrm{pc}}t}.
    }
    \label{eq:AU_global_phase_avalanche}
\end{equation}

The exponential form therefore represents the mean-field evolution of a
self-reproducing live-boundary network with an approximately stationary net
fractional growth rate.


% ---------------------------------------------------------------------
\subsection{Relation Between the Phase-Avalanche Rate and Vacuum Relaxation}
% ---------------------------------------------------------------------

The structural Phase-Avalanche Rate
\(\Gamma_{\mathrm{pc}}\)
and the coherent photon Energy-Relaxation Rate
\(\Gamma_{\mathrm D}\)
describe distinct projections of the Active Vacuum.

The photon branch obeys
\begin{equation}
    \frac{
        \mathrm d\ln E_\gamma
    }{
        \mathrm dt
    }
    =
    -\Gamma_{\mathrm D},
\end{equation}
whereas the global State II boundary network obeys
\begin{equation}
    \frac{
        \mathrm d\ln M_{\mathrm p}
    }{
        \mathrm dt
    }
    =
    +\Gamma_{\mathrm{pc}}.
\end{equation}

Define the optical--structural rate ratio
\begin{equation}
    \boxed{
    \mathcal C_{\mathrm{pcD}}
    \equiv
    \frac{
        \Gamma_{\mathrm{pc}}
    }{
        \Gamma_{\mathrm D}
    }.
    }
    \label{eq:AU_pcD_ratio}
\end{equation}

A common slow Active-Vacuum response may project differently onto the photon
scale coordinate and the global phase-boundary network. The single-rate TFT
specialization is
\begin{equation}
    \boxed{
    \mathcal C_{\mathrm{pcD}}
    =
    1
    \quad\Longrightarrow\quad
    \Gamma_{\mathrm{pc}}
    =
    \Gamma_{\mathrm D}.
    }
    \label{eq:AU_single_rate_specialization}
\end{equation}

No direct conversion of coherently relaxed photon energy into Black Hole mass
is implied by
Equation~\eqref{eq:AU_single_rate_specialization}. The two channels remain
mechanically distinct.

For the reference numerical chronology below, adopt the unit-reference
single-rate specialization
\begin{equation}
    \boxed{
    \Gamma_{\mathrm{pc}}^{(0)}
    =
    \Gamma_{\mathrm D}^{(0)}
    =
    \frac{9}{16}
    \frac{
        c_{\mathrm E}
    }{
        r_{\mathrm D}
    }
    \left(
        \frac{
            \ell_{\mathrm S}
        }{
            r_{\mathrm D}
        }
    \right)^6
    \simeq
    2.3591\times10^{-18}
    \ {\rm s^{-1}}.
    }
    \label{eq:AU_gamma_pc_reference}
\end{equation}

Equation~\eqref{eq:AU_gamma_pc_reference}
provides the reference synchronized timescale. The structural chronology
itself remains governed generally by
\(\Gamma_{\mathrm{pc}}\).

The corresponding reference e-folding time is
\begin{equation}
    \frac{
        1
    }{
        \Gamma_{\mathrm{pc}}^{(0)}
    }
    \simeq
    13.432
    \ {\rm Gyr}.
    \label{eq:AU_pc_efold_time}
\end{equation}

For the kinematic-length representation,
\begin{equation}
    x
    =
    c_{\mathrm E}t.
\end{equation}

This is a scale representation of the elapsed structural clock and is not
identified with ordinary radial material motion of the full Structural
Reservoir.

Equation~\eqref{eq:AU_global_phase_avalanche} may then be written
\begin{equation}
    M_{\mathrm p}(x)
    =
    M_{\mathrm p}(0)
    \exp
    \left(
        \frac{
            \Gamma_{\mathrm{pc}}x
        }{
            c_{\mathrm E}
        }
    \right).
    \label{eq:AU_phase_avalanche_length}
\end{equation}


% ---------------------------------------------------------------------
\subsection{Variable Definitions and Initialization Boundary}
% ---------------------------------------------------------------------

\begin{itemize}

    \item \textbf{\(M_{\mathrm p}(t)\) (Aggregate Percolating Mass):}
    The total inertial mass phase-segregated onto the distributed live FCC/SC
    boundaries participating in the global State II network.

    \item \textbf{\(M_{\mathrm p}(0)\) (Primordial Aggregate Boundary Mass):}
    The effective aggregate live-boundary mass established during the short
    primordial State II nucleation interval following Catastrophic
    Decompression.

    \item \textbf{\(M_{\mathrm{D,G}}\) (Micro-Black Hole Seed):}
    Approximately \(10.033~{\rm kg}\). This is the minimum localized
    kinematic mass threshold required to nucleate one stable State II phase
    domain. Therefore,
    \begin{equation}
        M_{\mathrm p}(0)
        \geq
        M_{\mathrm{D,G}}.
    \end{equation}

    \item \textbf{\(\Gamma_{\mathrm{pc}}\) (Phase-Avalanche Rate):}
    The net fractional growth rate of the aggregate live FCC/SC
    phase-boundary network,
    \begin{equation}
        \Gamma_{\mathrm{pc}}
        =
        \frac{
            \dot M_{\mathrm p}
        }{
            M_{\mathrm p}
        }.
    \end{equation}

    \item \textbf{\(\Gamma_{\mathrm D}\) (Vacuum Relaxation Rate):}
    The fractional coherent-photon scale-relaxation rate. Under the
    single-rate Active-Vacuum specialization,
    \begin{equation}
        \Gamma_{\mathrm{pc}}
        =
        \Gamma_{\mathrm D}.
    \end{equation}

\end{itemize}

\textbf{The Initialization Boundary (\(t\approx0\))}

Because State II domains are lattice phase transitions rather than atomic
structures, their nucleation does not require the formation of neutral matter.

In the immediate aftermath of Catastrophic Decompression, overlapping
longitudinal acoustic shockwaves can produce regions of extreme localized
compression.

Wherever the active FCC vacuum reaches the saturation condition
\begin{equation}
    \alpha_{\mathrm G}
    \rightarrow
    1,
\end{equation}
the local Soft Floor becomes phase locked,
\begin{equation}
    \mathrm{FCC}
    \rightarrow
    \mathrm{SC},
\end{equation}
creating primordial State II nuclei.

A potentially large population of primordial nuclei may therefore form during
the short post-rupture transient. Primordial nucleation occurring over a time
interval negligible relative to the subsequent structural e-folding time is
absorbed into the effective initialization
\(M_{\mathrm p}(0)\).

The minimum possible initialization is the single-seed limit,
\begin{equation}
    M_{\mathrm p}(0)
    =
    M_{\mathrm{D,G}},
\end{equation}
while any distributed primordial population satisfies
\begin{equation}
    M_{\mathrm p}(0)
    >
    M_{\mathrm{D,G}}.
\end{equation}

The phase-avalanche clock therefore begins with restoration of the active FCC
vacuum and precedes the later lattice-annealing epoch required for stable
neutral hydrogen.


% ---------------------------------------------------------------------
\subsection{Integration of the Macroscopic Lifespan
\texorpdfstring{$T_{\mathrm{cyc}}$}{Tcyc}}
% ---------------------------------------------------------------------

As derived in
\hyperref[sec:reservoir_mass]{Section~\ref*{sec:reservoir_mass}},
the total structural reservoir mass is
\begin{equation}
    M_{\mathrm{Res}}
    \approx
    1.6707\times10^{79}
    \ {\rm kg}.
\end{equation}

Under the equal-capacity percolation closure, the cycle terminates when the
aggregate live-boundary mass reaches
\begin{equation}
    M_{\mathrm{trg}}
    \approx
    8.6208\times10^{77}
    \ {\rm kg}.
\end{equation}

Applying
Equation~\eqref{eq:AU_global_phase_avalanche}
at the rupture time gives
\begin{equation}
    M_{\mathrm{trg}}
    =
    M_{\mathrm p}(0)
    e^{\Gamma_{\mathrm{pc}}T_{\mathrm{cyc}}}.
\end{equation}

Therefore,
\begin{equation}
    \boxed{
    \Gamma_{\mathrm{pc}}
    T_{\mathrm{cyc}}
    =
    \ln
    \left(
        \frac{
            M_{\mathrm{trg}}
        }{
            M_{\mathrm p}(0)
        }
    \right).
    }
    \label{eq:AU_cycle_depth_general}
\end{equation}

The structural cycle time is consequently
\begin{equation}
    \boxed{
    T_{\mathrm{cyc}}
    =
    \frac{
        1
    }{
        \Gamma_{\mathrm{pc}}
    }
    \ln
    \left(
        \frac{
            M_{\mathrm{trg}}
        }{
            M_{\mathrm p}(0)
        }
    \right).
    }
    \label{eq:AU_cycle_time_general}
\end{equation}

Equation~\eqref{eq:AU_cycle_depth_general}
separates the two ingredients of the chronology:

\begin{equation}
    \boxed{
    \text{percolation geometry fixes the required structural depth,}
    \qquad
    \Gamma_{\mathrm{pc}}^{-1}
    \text{ fixes the duration of one e-fold}.
    }
\end{equation}

Because
\begin{equation}
    M_{\mathrm p}(0)
    \geq
    M_{\mathrm{D,G}},
\end{equation}
the minimum-seed initialization produces the maximum possible cycle duration
for a fixed
\(\Gamma_{\mathrm{pc}}\):
\begin{equation}
    T_{\mathrm{cyc}}
    \leq
    \frac{
        1
    }{
        \Gamma_{\mathrm{pc}}
    }
    \ln
    \left(
        \frac{
            M_{\mathrm{trg}}
        }{
            M_{\mathrm{D,G}}
        }
    \right).
\end{equation}

For the minimum-seed branch,
\begin{equation}
    M_{\mathrm p}(0)
    =
    M_{\mathrm{D,G}}
    \approx
    10.033
    \ {\rm kg},
\end{equation}
so
\begin{equation}
    \frac{
        M_{\mathrm{trg}}
    }{
        M_{\mathrm{D,G}}
    }
    \approx
    8.5924\times10^{76}.
\end{equation}

Thus,
\begin{equation}
    \boxed{
    \Gamma_{\mathrm{pc}}
    T_{\mathrm{cyc}}^{(\mathrm{max})}
    =
    \ln
    \left(
        8.5924\times10^{76}
    \right)
    \approx
    177.15.
    }
    \label{eq:AU_cycle_efolds}
\end{equation}

The minimum-seed universe therefore requires approximately
\(177.15\)
structural e-folds to evolve from one minimum stable State II seed to the
global rupture threshold under the exponential mean-field branch.

Using the reference single-rate value
\begin{equation}
    \Gamma_{\mathrm{pc}}^{(0)}
    \simeq
    2.3591\times10^{-18}
    \ {\rm s^{-1}},
\end{equation}
gives
\begin{equation}
    \boxed{
    T_{\mathrm{cyc}}^{(0,\mathrm{max})}
    \approx
    2.3795
    \ {\rm trillion\ years}.
    }
    \label{eq:AU_cycle_time_reference}
\end{equation}

This is the maximum cycle duration of the minimum-seed reference branch. A
larger primordial aggregate
\(M_{\mathrm p}(0)\)
shortens the cycle logarithmically but does not alter the rupture mass.

The corresponding kinematic cycle length is
\begin{equation}
    L_{\mathrm{max}}
    \equiv
    c_{\mathrm E}
    T_{\mathrm{cyc}},
\end{equation}
so generally
\begin{equation}
    \boxed{
    L_{\mathrm{max}}
    =
    \frac{
        c_{\mathrm E}
    }{
        \Gamma_{\mathrm{pc}}
    }
    \ln
    \left(
        \frac{
            M_{\mathrm{trg}}
        }{
            M_{\mathrm p}(0)
        }
    \right).
    }
    \label{eq:AU_cycle_length_general}
\end{equation}

Under the single-rate specialization
\begin{equation}
    \Gamma_{\mathrm{pc}}
    =
    \Gamma_{\mathrm D},
\end{equation}
the structural e-fold length becomes
\begin{equation}
    \frac{
        c_{\mathrm E}
    }{
        \Gamma_{\mathrm{pc}}
    }
    =
    \frac{
        c_{\mathrm E}
    }{
        \Gamma_{\mathrm D}
    }
    =
    R_{\mathrm U},
\end{equation}
and therefore
\begin{equation}
    L_{\mathrm{max}}
    =
    R_{\mathrm U}
    \ln
    \left(
        \frac{
            M_{\mathrm{trg}}
        }{
            M_{\mathrm p}(0)
        }
    \right).
\end{equation}

For the minimum-seed reference normalization,
\begin{equation}
    \boxed{
    L_{\mathrm{max}}^{(0)}
    \approx
    2.25\times10^{28}
    \ {\rm m}.
    }
\end{equation}

At the endpoint,
\begin{equation}
    x
    =
    L_{\mathrm{max}}
    \qquad\Longleftrightarrow\qquad
    t
    =
    T_{\mathrm{cyc}}
    \qquad\Longleftrightarrow\qquad
    M_{\mathrm p}
    =
    M_{\mathrm{trg}}.
\end{equation}

Thus \(L_{\mathrm{max}}\) is the kinematic-length representation of the
complete structural cycle. It is not an independent material or causal
boundary of the substrate.


% ---------------------------------------------------------------------
\subsection{Integration of the Current Epoch
\texorpdfstring{$t_{\mathrm{cur}}$}{tcur}}
% ---------------------------------------------------------------------

To estimate the current elapsed phase time, the Global Phase Avalanche is
evaluated against the present aggregate State II boundary mass
\(M_{\mathrm{cur}}\).

Adopting the current observational Black Hole mass inventory over the
reference observable domain as
\begin{equation}
    M_{\mathrm{obs}}
    \approx
    1.000\times10^{50}
    \ {\rm kg},
\end{equation}
TFT presently extrapolates this inventory across the full Structural Reservoir
using the volumetric reservoir multiplier
\begin{equation}
    9.7957\times10^{25},
\end{equation}
under the statistical-homogeneity closure that the observed State II
abundance is representative of the corresponding larger structural domain.

This gives
\begin{equation}
    M_{\mathrm{cur}}
    =
    \left(
        1.000\times10^{50}
        \ {\rm kg}
    \right)
    \left(
        9.7957\times10^{25}
    \right)
    =
    9.7957\times10^{75}
    \ {\rm kg}.
\end{equation}

The Global Phase Avalanche gives
\begin{equation}
    M_{\mathrm{cur}}
    =
    M_{\mathrm p}(0)
    e^{\Gamma_{\mathrm{pc}}t_{\mathrm{cur}}},
\end{equation}
and therefore
\begin{equation}
    \boxed{
    t_{\mathrm{cur}}
    =
    \frac{
        1
    }{
        \Gamma_{\mathrm{pc}}
    }
    \ln
    \left(
        \frac{
            M_{\mathrm{cur}}
        }{
            M_{\mathrm p}(0)
        }
    \right).
    }
    \label{eq:AU_current_time_general}
\end{equation}

For the minimum-seed normalization,
\begin{equation}
    \frac{
        M_{\mathrm{cur}}
    }{
        M_{\mathrm{D,G}}
    }
    \approx
    9.7635\times10^{74},
\end{equation}
so
\begin{equation}
    \boxed{
    \Gamma_{\mathrm{pc}}
    t_{\mathrm{cur}}
    \approx
    172.67.
    }
    \label{eq:AU_current_efolds}
\end{equation}

Using the reference single-rate value
\(\Gamma_{\mathrm{pc}}^{(0)}\),
\begin{equation}
    \boxed{
    t_{\mathrm{cur}}^{(0)}
    \approx
    2.3194
    \ {\rm trillion\ years}.
    }
    \label{eq:AU_current_time_reference}
\end{equation}

As with
\(T_{\mathrm{cyc}}\),
this absolute elapsed-time value corresponds to the minimum-seed normalization
and the adopted statistical-homogeneity extrapolation.

A larger primordial aggregate shifts
\(t_{\mathrm{cur}}\)
and
\(T_{\mathrm{cyc}}\)
downward by the same logarithmic initialization term.


% ---------------------------------------------------------------------
\subsection{The Expiration Timer and Remaining Lifespan}
% ---------------------------------------------------------------------

The differential time remaining before global percolation does not depend on
the primordial normalization
\(M_{\mathrm p}(0)\).

From
\begin{equation}
    M_{\mathrm{trg}}
    =
    M_{\mathrm p}(0)
    e^{\Gamma_{\mathrm{pc}}T_{\mathrm{cyc}}},
\end{equation}
and
\begin{equation}
    M_{\mathrm{cur}}
    =
    M_{\mathrm p}(0)
    e^{\Gamma_{\mathrm{pc}}t_{\mathrm{cur}}},
\end{equation}
their ratio is
\begin{equation}
    \frac{
        M_{\mathrm{trg}}
    }{
        M_{\mathrm{cur}}
    }
    =
    e^{
        \Gamma_{\mathrm{pc}}
        \left(
            T_{\mathrm{cyc}}
            -
            t_{\mathrm{cur}}
        \right)
    }.
\end{equation}

Define
\begin{equation}
    \Delta t
    \equiv
    T_{\mathrm{cyc}}
    -
    t_{\mathrm{cur}}.
\end{equation}

Then
\begin{equation}
    \boxed{
    \Delta t
    =
    \frac{
        1
    }{
        \Gamma_{\mathrm{pc}}
    }
    \ln
    \left(
        \frac{
            M_{\mathrm{trg}}
        }{
            M_{\mathrm{cur}}
        }
    \right).
    }
    \label{eq:AU_remaining_time_general}
\end{equation}

Using
\begin{equation}
    \frac{
        M_{\mathrm{trg}}
    }{
        M_{\mathrm{cur}}
    }
    =
    \frac{
        8.6208\times10^{77}
    }{
        9.7957\times10^{75}
    }
    \approx
    88.006,
\end{equation}
gives
\begin{equation}
    \boxed{
    \Gamma_{\mathrm{pc}}
    \Delta t
    =
    \ln(88.006)
    \approx
    4.4774.
    }
    \label{eq:AU_remaining_efolds}
\end{equation}

Thus the present reference state lies approximately
\(4.48\)
structural e-folds below the global percolation threshold.

Using
\begin{equation}
    \Gamma_{\mathrm{pc}}^{(0)}
    \simeq
    2.3591\times10^{-18}
    \ {\rm s^{-1}},
\end{equation}
the reference remaining time is
\begin{equation}
    \boxed{
    \Delta t^{(0)}
    \approx
    60.14
    \ {\rm billion\ years}.
    }
    \label{eq:AU_remaining_time_reference}
\end{equation}

The cancellation of
\(M_{\mathrm p}(0)\)
is exact within the exponential Global Phase Avalanche law.

The numerical value depends primarily on the reference value of
\(\Gamma_{\mathrm{pc}}\), the percolation-derived trigger mass, and the
statistical-homogeneity extrapolation used to estimate
\(M_{\mathrm{cur}}\).


% ---------------------------------------------------------------------
\subsection{Conditional Optical--Structural Clock Synchronization}
% ---------------------------------------------------------------------

The coherent optical relaxation clock and the structural Phase-Avalanche clock
may now be compared without identifying their rates a priori.

For a coherent transverse excitation propagating through the constant-rate,
weak-loading optical branch,
\begin{equation}
    \frac{
        E_{\mathrm{obs}}(t)
    }{
        E_{\mathrm{emit}}
    }
    =
    e^{-\Gamma_{\mathrm D}t}.
    \label{eq:AU_optical_clock}
\end{equation}

The Global Phase Avalanche gives
\begin{equation}
    \frac{
        M_{\mathrm p}(t)
    }{
        M_{\mathrm p}(0)
    }
    =
    e^{+\Gamma_{\mathrm{pc}}t}.
    \label{eq:AU_structural_clock}
\end{equation}

Eliminating \(t\) between
Equations~\eqref{eq:AU_optical_clock}
and~\eqref{eq:AU_structural_clock}
gives the general optical--structural clock relation
\begin{equation}
    \boxed{
    \frac{
        E_{\mathrm{obs}}(t)
    }{
        E_{\mathrm{emit}}
    }
    =
    \left[
        \frac{
            M_{\mathrm p}(t)
        }{
            M_{\mathrm p}(0)
        }
    \right]^{
        -\Gamma_{\mathrm D}/\Gamma_{\mathrm{pc}}
    }.
    }
    \label{eq:AU_general_clock_relation}
\end{equation}

Using
Equation~\eqref{eq:AU_pcD_ratio},
this may also be written
\begin{equation}
    \frac{
        E_{\mathrm{obs}}(t)
    }{
        E_{\mathrm{emit}}
    }
    =
    \left[
        \frac{
            M_{\mathrm p}(t)
        }{
            M_{\mathrm p}(0)
        }
    \right]^{-1/\mathcal C_{\mathrm{pcD}}}.
\end{equation}

The single-rate synchronization specialization
\begin{equation}
    \Gamma_{\mathrm{pc}}
    =
    \Gamma_{\mathrm D}
\end{equation}
reduces
Equation~\eqref{eq:AU_general_clock_relation}
to the reciprocal identity
\begin{equation}
    \boxed{
    \frac{
        E_{\mathrm{obs}}(t)
    }{
        E_{\mathrm{emit}}
    }
    =
    \frac{
        M_{\mathrm p}(0)
    }{
        M_{\mathrm p}(t)
    }.
    }
    \label{eq:AU_reciprocal_clock}
\end{equation}

Thus reciprocal optical and structural exponentials are not assumed at the
outset. They emerge as the special case in which the two Active-Vacuum
projections possess the same fractional rate.

For the optical branch,
\begin{equation}
    R_{\mathrm U}
    =
    \frac{
        c_{\mathrm E}
    }{
        \Gamma_{\mathrm D}
    }
\end{equation}
is the characteristic one-\(e\)-fold coherent-relaxation length, with
\begin{equation}
    E_{\mathrm{obs}}(R_{\mathrm U})
    =
    \frac{
        E_{\mathrm{emit}}
    }{
        e
    }.
\end{equation}

The scale
\(R_{\mathrm U}\)
is not a universal stopping radius. A physical terminal coherence-loss
condition, if one exists, is determined separately by photon-branch
stability and is not required for the formal clock comparison below.

The optical exponential may therefore be mathematically continued to the
structural endpoint
\(T_{\mathrm{cyc}}\).

At the Big Rupture,
\begin{equation}
    \frac{
        M_{\mathrm{trg}}
    }{
        M_{\mathrm p}(0)
    }
    =
    e^{\Gamma_{\mathrm{pc}}T_{\mathrm{cyc}}},
\end{equation}
so the formally continued optical factor at the same elapsed structural time
is
\begin{equation}
    \boxed{
    \frac{
        E_{\mathrm{obs}}^{(\mathrm{formal})}
        (T_{\mathrm{cyc}})
    }{
        E_{\mathrm{emit}}
    }
    =
    \left(
        \frac{
            M_{\mathrm{trg}}
        }{
            M_{\mathrm p}(0)
        }
    \right)^{
        -\Gamma_{\mathrm D}/\Gamma_{\mathrm{pc}}
    }.
    }
    \label{eq:AU_endpoint_clock_general}
\end{equation}

Under the single-rate specialization,
\begin{equation}
    \boxed{
    \frac{
        E_{\mathrm{obs}}^{(\mathrm{formal})}
        (T_{\mathrm{cyc}})
    }{
        E_{\mathrm{emit}}
    }
    =
    \frac{
        M_{\mathrm p}(0)
    }{
        M_{\mathrm{trg}}
    }.
    }
    \label{eq:AU_endpoint_clock_single_rate}
\end{equation}

For the minimum-seed normalization,
\begin{equation}
    M_{\mathrm p}(0)
    =
    M_{\mathrm{D,G}},
\end{equation}
this becomes
\begin{equation}
    \frac{
        E_{\mathrm{obs}}^{(\mathrm{formal})}
        (T_{\mathrm{cyc}})
    }{
        E_{\mathrm{emit}}
    }
    =
    \frac{
        M_{\mathrm{D,G}}
    }{
        M_{\mathrm{trg}}
    }.
\end{equation}

If the emitted excitation is additionally normalized to the established
maximum-energy scale,
\begin{equation}
    E_{\mathrm{emit}}
    =
    E_{\mathrm{max}}
    \approx
    13.331
    \ {\rm J},
\end{equation}
then
\begin{equation}
    E_{\mathrm{obs}}^{(\mathrm{formal})}
    (T_{\mathrm{cyc}})
    =
    E_{\mathrm{max}}
    \left(
        \frac{
            M_{\mathrm{D,G}}
        }{
            M_{\mathrm{trg}}
        }
    \right).
\end{equation}

Using
\begin{equation}
    M_{\mathrm{D,G}}
    \approx
    10.033
    \ {\rm kg},
\end{equation}
and
\begin{equation}
    M_{\mathrm{trg}}
    \approx
    8.6208\times10^{77}
    \ {\rm kg},
\end{equation}
gives
\begin{equation}
    E_{\mathrm{obs}}^{(\mathrm{formal})}
    (T_{\mathrm{cyc}})
    \approx
    1.5515\times10^{-76}
    \ {\rm J}.
\end{equation}

This value is the formal continuation of the coherent optical relaxation
factor to the structural-cycle endpoint. It does not require an individual
photon to remain on the coherent branch throughout the complete structural
cycle.

The reciprocal single-rate identity likewise does not imply that coherently
relaxed photon energy is converted directly into State II boundary mass.

The optical channel and the structural phase-avalanche channel remain
mechanically distinct. Their synchronization expresses equality of their
fractional Active-Vacuum rates under
Equation~\eqref{eq:AU_single_rate_specialization}.

The synchronized structural endpoint is
\begin{equation}
    x
    =
    L_{\mathrm{max}}
    \qquad\Longleftrightarrow\qquad
    t
    =
    T_{\mathrm{cyc}}
    \qquad\Longleftrightarrow\qquad
    M_{\mathrm p}
    =
    M_{\mathrm{trg}}.
\end{equation}

The resulting hierarchy is therefore
\begin{equation}
    \boxed{
    \text{live-boundary conservation}
    \rightarrow
    \Gamma_{\mathrm{pc}}
    \rightarrow
    \text{Global Phase Avalanche}
    \rightarrow
    \text{percolation trigger}
    \rightarrow
    T_{\mathrm{cyc}}.
    }
\end{equation}

Under the additional single-rate Active-Vacuum specialization,
\begin{equation}
    \Gamma_{\mathrm{pc}}
    =
    \Gamma_{\mathrm D},
\end{equation}
the structural and optical clocks become reciprocal projections of one common
macroscopic relaxation timescale.